NHỚ MÃI CÁC ANH-NHỮNG NGƯỜI XÂY NỀN MÓNG TIN HỌC VIỆT NAM.

Sau 1975, với sự thành lập Viện Khoa học tính toán và Điều khiển (Viện Khoa học Việt Nam) do GS Phan Đình Diệu làm Viện trưởng và Cục Máy tính điện tử (Ủy ban khoa học và kỹ thuật nhà nước) do GS Trần lưu Chương làm Cục trưởng ,Tin học Việt Nam đã có bước phát triển mới. Các Anh đã  lãnh đạo 2 đơn vị hợp tác, triển khai thành công xuất sắc một khối lượng lớn công việc trong giai đoạn đầu của nền khoa học và công nghệ non trẻ này, đồng thời xây dựng nền móng choTin học Việt Nam:

1-Những Nghị quyết, Nghị định, Chính sách về các hoạt động nghiên cứu, ứng dụng Máy tính điện tử và phát triển Tin hoc Việt Nam trong thập niên 80,90..đều ghi dấu ấn trí tuệ sâu sắc và tâm huyết của các Anh,mở đường cho Thế hệ sau phát triển ..

2-Đề xuất và trực tiếp tham gia tổ chức thực hiện nhiều chương trình, dự án trọng điểm với tầm nhìn, hoài bão cháy bỏng về CNTT Việt Nam ngang tầm thế giới ( Điển hình là dự án thiết kế chế tạo máy vi tính VT80,81).

3-Với cương vị lãnh đạo cùng tài năng uyên bác, đức độ của mình,các Anh đã kết nối được nhiều tổ chức và nhiều nhà khoa học ,chuyên gia quốc tế  xuất sắc giúp đỡ rất  hiệu quả  cho nền Tin học Việt Nam  non trẻ trong nhiều năm bị cấm vận!

4-Mặc dù làm công tác quản lý ,các Anh vẫn say mê nghiên cứu và tham gia giảng dạy,đào tạo lớp trẻ. Các Anh đã tổ chức được nhiều chương trình hợp tác quốc tế  đào tạo,bồi dưỡng trình độ cao  cho nhiều  thế hệ các nhà chuyên môn và quản lý đang  phát huy tài năng rất tốt trong nhiều trường Đại học, Viện, Ngành..trong cả nước. 

Tôi cũng rất vinh dự được tham gia và giúp một số việc nhỏ cho các Anh và cũng học hỏi được rất nhiều. Sinh thời,các Anh rất quý mến và thân thiết với nhau.Đặc biệt, tôi cũng được các Anh dành cho một ân tình:
Năm 1982,Anh Diệu đã giới thiệu tôi ứng viên Cộng tác viên khoa học bậc cao vào Phòng thí nghiệm quốc tế về Trí tuệ nhân tạo ( Bratislava -Tiệp khắc),Anh Chương đã lo thủ tục và vận động tổ chức để tôi được đi..Xin thắp nén tâm hương nhớ mãi các Anh.